\section{Projective schemes and computational algebraic geometry}

\subsection{The standard affine charts of \texorpdfstring{$\Proj$}{Proj}}

Let
\[
 S=R[X_0,\ldots,X_n]/I
\]
be a finitely presented graded ring, with $I$ given by finitely many homogeneous relations.  The
standard chart $D_+(X_i)$ is represented algebraically by the degree-zero part of the localization
$S_{X_i}$.  For a projective subscheme of $\mathbf P^n_R$, this chart has the familiar finite
presentation obtained by setting $X_i=1$ and introducing affine coordinates
\[
                    x_{ij}=X_j/X_i\qquad(j\neq i).
\]
Each homogeneous relation is dehomogenized with respect to $X_i$.

The overlap of charts $i$ and $j$ is the principal open in chart $i$ where $x_{ij}$ is invertible.
The transition map is
\[
 x_{jk}\longmapsto x_{ik}/x_{ij},
 \qquad
 x_{ji}\longmapsto 1/x_{ij}.
\]
These formulas are maps of localization presentations.  Their inverse and triple-cocycle equations
are finite rational-function identities, hence localization equality certificates.

\begin{proposition}[Finite atlas for a projective presentation]\label{prop:proj-atlas}
For every fixed finite homogeneous presentation of $S$, the $n+1$ standard charts, their principal
overlaps, transition maps, inverse witnesses, and triple cocycles are primitive-recursively
constructible.  They form a strict certified affine atlas.
\end{proposition}

\begin{proof}
Dehomogenization is a bounded substitution on each finite polynomial code.  The transition maps are
fixed monomial substitutions into localizations.  Applying a homogeneous relation and clearing the
explicit denominator gives its target relation certificate.  The inverse and cocycle equations are
monomial exponent identities; a bounded coefficient checker verifies them.  The number of charts
and overlaps is finite and bounded by polynomial expressions in $n$.
\end{proof}

Thus a large class of projective geometry can be developed without defining $\Proj$ internally as a
set of homogeneous prime ideals.

\subsection{Homogeneous coordinates and points}

For a field or domain with a suitable exact presentation, a projective point may be represented by a
nonzero tuple $(x_0,\ldots,x_n)$ modulo common scaling.  The existing Diophantine geometry packet
uses primitive representatives and height.  The atlas representation and homogeneous-coordinate
representation are complementary:

\begin{itemize}
\item homogeneous coordinates are compact for point enumeration and height arguments;
\item affine charts are natural for local rings, functions, sheaves, Jacobians, and descent;
\item a point with $x_i$ invertible maps to chart $i$ by the ratios $x_j/x_i$;
\item two chart descriptions agree by the certified transition maps.
\end{itemize}

A theorem may use a finite relational choice of a nonzero coordinate.  If its final output is
projectively invariant, the finite-use framework eliminates dependence on the selected chart.
Alternatively, a primitive representative over $\Z$ can choose the least nonzero coordinate
deterministically.

\subsection{Projective morphisms}

A tuple of homogeneous forms $F_0,\ldots,F_m$ of one common degree defines a projective morphism on
the locus where they do not vanish simultaneously.  A finite certificate for a global morphism on a
projective subscheme consists of:

\begin{enumerate}[label=(\alph*)]
\item relation witnesses showing that the image satisfies the target homogeneous equations;
\item a finite principal-cover certificate that the $F_j$ generate the unit ideal on each source
      chart, or equivalent local nonvanishing certificates;
\item chartwise rational formulas and overlap commutativity witnesses.
\end{enumerate}

Composition substitutes homogeneous forms and multiplies degrees.  The relation and base-locus
certificates transform by finite polynomial identities.  Height upper bounds are obtained from
coefficient norms exactly as in the Diophantine packet; lower bounds still require gcd/resultant or
base-locus information and are not produced by the atlas formalism alone.

\subsection{Jacobian and smoothness certificates}

For a closed affine presentation $A=R[X]/(g_1,\ldots,g_s)$, the Jacobian matrix is a finite matrix of
formal partial derivatives.  A certificate that a chosen $d\times d$ minor is invertible on a
principal open consists of a localization inverse witness.  A finite list of such minors and a
unit-ideal cover certificate proves that the required rank condition holds everywhere on the chosen
atlas.

This gives an iPRA-friendly sufficient interface for smoothness in standard complete-intersection
situations.  It does not assert a uniform algorithm deciding smoothness of arbitrary presentations.
The input may carry a Gröbner/elimination certificate producing the relevant Jacobian ideal, or the
minor cover may simply be part of the theorem hypotheses.

\subsection{Divisors and line bundles}

A Cartier divisor on a finite atlas is represented by local nonzerodivisor fractions $g_i$ and unit
transition witnesses $g_i/g_j\in\cO^\times(U_i\cap U_j)$.  The associated line bundle uses those
units as rank-one transition matrices.  Tensor product adds divisor data and multiplies transition
units; dualization inverts them on overlaps.

For projective space, the hyperplane divisor gives $\cO(1)$, and $\cO(d)$ has transition factor
$(X_j/X_i)^d$.  All integer powers are finite localization terms.  The projective-line and
projective-space computations below therefore require no abstract Picard group.

\subsection{Hilbert functions and finite graded certificates}

A homogeneous Gröbner certificate can provide a finite monomial ideal of leading terms.  For a fixed
degree $D$, the standard monomials of degree $D$ are a bounded finite set, and coefficient reduction
proves a spanning or basis statement for the degree-$D$ quotient.  Hence the Hilbert function at a
specified degree is a finite iPRA observation.

A claim about the eventual Hilbert polynomial needs an additional stabilization certificate: a
bound $D_0$, a polynomial $P$, and finite algebraic data proving the relevant recurrence or graded
resolution for all $D\geq D_0$.  The framework checks such a certificate; it does not infer an
unbounded stabilization theorem from a finite table.

\subsection{Elimination, images, and fibers}

For a presented morphism, an elimination certificate can describe equations for a scheme-theoretic
image or a fiber.  Soundness is again finite: substitute the morphism tuple and verify ideal
membership.  Completeness of the proposed elimination ideal requires a Gröbner basis, resultant, or
other finite certificate appropriate to the family.

This division of labor is useful in iPRA:

\begin{itemize}
\item polynomial substitution and certificate verification are generic;
\item the domain-specific algorithm produces the elimination certificate;
\item the geometric theorem consumes only the certified presentation of the resulting image or
      fiber.
\end{itemize}

\subsection{What part of algebraic geometry is thereby accessible?}

The finite-atlas language directly supports substantial fragments of:

\begin{itemize}
\item affine and projective varieties/schemes of finite presentation;
\item regular and rational maps on certified domains;
\item finite morphisms and finite module algebras;
\item closed immersions, fiber products, and base change;
\item vector bundles, line bundles, and finite \v{C}ech complexes;
\item Jacobian calculations and chartwise smoothness certificates;
\item exact finite-field point counting by bounded enumeration;
\item rational/integral point searches once an external height bound is supplied;
\item graded pieces, Hilbert-function observations, and certificate-relative elimination.
\end{itemize}

The framework is intentionally less ambitious about arbitrary non-quasi-compact schemes, unbounded
sites, derived categories as completed universes, and existence theorems whose proofs provide no
primitive-recursive quantitative data.  Those topics can still have finite theorem-specific
observations, but they are not forced into a single universal code.
